% ===== Ядро списку джерел. Ключі є СПІЛЬНИМИ для всіх розділів. =====
% Розділи НЕ дублюють ці позиції у власних файлах джерел, а лише посилаються \cite{ключ}.

% --- Законодавство України ---
\bibitem{zu-info} Про інформацію: Закон України від 02.10.1992 № 2657-XII. URL: \url{https://zakon.rada.gov.ua/laws/show/2657-12}.

\bibitem{zu-zi} Про захист інформації в інформаційно-комунікаційних системах: Закон України від 05.07.1994 № 80/94-ВР. URL: \url{https://zakon.rada.gov.ua/laws/show/80/94-вр}.

\bibitem{zu-cyber} Про основні засади забезпечення кібербезпеки України: Закон України від 05.10.2017 № 2163-VIII (редакція від 19.10.2025). URL: \url{https://zakon.rada.gov.ua/laws/show/2163-19}.

\bibitem{zu-ki} Про критичну інфраструктуру: Закон України від 16.11.2021 № 1882-IX. URL: \url{https://zakon.rada.gov.ua/laws/show/1882-20}.

\bibitem{zu-pd} Про захист персональних даних: Закон України від 01.06.2010 № 2297-VI. URL: \url{https://zakon.rada.gov.ua/laws/show/2297-17}.

\bibitem{kk-ukr} Кримінальний кодекс України від 05.04.2001 № 2341-III (розділ XVI «Кримінальні правопорушення у сфері використання ЕОМ, систем та мереж електрозв'язку»). URL: \url{https://zakon.rada.gov.ua/laws/show/2341-14}.

\bibitem{zu-nb} Про національну безпеку України: Закон України від 21.06.2018 № 2469-VIII. URL: \url{https://zakon.rada.gov.ua/laws/show/2469-19}.

\bibitem{ukaz-447} Про Стратегію кібербезпеки України: Указ Президента України від 26.08.2021 № 447/2021. URL: \url{https://zakon.rada.gov.ua/laws/show/447/2021}.

\bibitem{kmu-518} Загальні вимоги до кіберзахисту об'єктів критичної інфраструктури: затв. постановою Кабінету Міністрів України від 19.06.2019 № 518. URL: \url{https://zakon.rada.gov.ua/laws/show/518-2019-п}.

% --- Регулятори та CERT ---
\bibitem{cip} Державна служба спеціального зв'язку та захисту інформації України: офіційний вебсайт. URL: \url{https://cip.gov.ua/}.

\bibitem{cert-ua} CERT-UA: Урядова команда реагування на комп'ютерні надзвичайні події України. URL: \url{https://cert.gov.ua/}.

\bibitem{scpc} Державний центр кіберзахисту Держспецзв'язку України. URL: \url{https://scpc.gov.ua/}.

% --- Нормативні документи ТЗІ та стандарти України ---
\bibitem{ndtzi-1-1-003} НД ТЗІ 1.1-003-99. Термінологія в галузі захисту інформації в комп'ютерних системах від несанкціонованого доступу. Київ: ДСТСЗІ СБ України, 1999. URL: \url{https://cip.gov.ua/ua/news/normativni-dokumenti-z-pitan-tzi}.

\bibitem{ndtzi-2-5-004} НД ТЗІ 2.5-004-99. Критерії оцінки захищеності інформації в комп'ютерних системах від несанкціонованого доступу. Київ: ДСТСЗІ СБ України, 1999.

\bibitem{ndtzi-3-7-003} НД ТЗІ 3.7-003-05. Порядок проведення робіт із створення комплексної системи захисту інформації в інформаційно-телекомунікаційній системі. Київ: ДСТСЗІ СБ України, 2005.

\bibitem{dstu-8302} ДСТУ 8302:2015. Інформація та документація. Бібліографічне посилання. Загальні положення та правила складання. Київ: ДП «УкрНДНЦ», 2016.

\bibitem{dstu-27001} ДСТУ EN ISO/IEC 27001:2023. Інформаційна безпека, кібербезпека та захист приватності. Системи керування інформаційною безпекою. Вимоги. Київ: ДП «УкрНДНЦ», 2023.

\bibitem{dstu-27002} ДСТУ EN ISO/IEC 27002:2023. Інформаційна безпека, кібербезпека та захист приватності. Заходи інформаційної безпеки. Київ: ДП «УкрНДНЦ», 2023.

\bibitem{dstu-7624} ДСТУ 7624:2014. Інформаційні технології. Криптографічний захист інформації. Алгоритм симетричного блокового перетворення. Київ: Мінекономрозвитку України, 2015.

\bibitem{dstu-7564} ДСТУ 7564:2014. Інформаційні технології. Криптографічний захист інформації. Функція гешування. Київ: Мінекономрозвитку України, 2015.

% --- Міжнародні стандарти ISO/IEC ---
\bibitem{iso-27035} ISO/IEC 27035-1:2023. Information technology --- Information security incident management --- Part~1: Principles and process. Geneva: ISO, 2023. URL: \url{https://www.iso.org/standard/78973.html}.

\bibitem{iso-27037} ISO/IEC 27037:2012. Guidelines for identification, collection, acquisition and preservation of digital evidence. Geneva: ISO, 2012. URL: \url{https://www.iso.org/standard/44381.html}.

\bibitem{iso-27005} ISO/IEC 27005:2022. Information security, cybersecurity and privacy protection --- Guidance on managing information security risks. Geneva: ISO, 2022. URL: \url{https://www.iso.org/standard/80585.html}.

\bibitem{iso-31000} ISO 31000:2018. Risk management --- Guidelines. Geneva: ISO, 2018. URL: \url{https://www.iso.org/standard/65694.html}.

\bibitem{iso-22301} ISO 22301:2019. Security and resilience --- Business continuity management systems --- Requirements. Geneva: ISO, 2019.

\bibitem{iec-62443} IEC 62443-3-3:2013. Industrial communication networks --- Network and system security --- Part~3-3: System security requirements and security levels. Geneva: IEC, 2013.

% --- NIST ---
\bibitem{nist-csf2} NIST. The NIST Cybersecurity Framework (CSF) 2.0. NIST CSWP 29. Gaithersburg, MD, 2024. DOI: 10.6028/NIST.CSWP.29.

\bibitem{nist-83} Souppaya M., Scarfone K. Guide to Malware Incident Prevention and Handling for Desktops and Laptops. NIST SP 800-83, Rev.~1. Gaithersburg, MD: NIST, 2013. DOI: 10.6028/NIST.SP.800-83r1.

\bibitem{nist-61} Cichonski P., Millar T., Grance T., Scarfone K. Computer Security Incident Handling Guide. NIST SP 800-61, Rev.~2. Gaithersburg, MD: NIST, 2012. DOI: 10.6028/NIST.SP.800-61r2.

\bibitem{nist-86} Kent K., Chevalier S., Grance T., Dang H. Guide to Integrating Forensic Techniques into Incident Response. NIST SP 800-86. Gaithersburg, MD: NIST, 2006. DOI: 10.6028/NIST.SP.800-86.

\bibitem{nist-137} Dempsey K. et al. Information Security Continuous Monitoring (ISCM) for Federal Information Systems and Organizations. NIST SP 800-137. Gaithersburg, MD: NIST, 2011. DOI: 10.6028/NIST.SP.800-137.

\bibitem{nist-53} NIST. Security and Privacy Controls for Information Systems and Organizations. NIST SP 800-53, Rev.~5. Gaithersburg, MD: NIST, 2020. DOI: 10.6028/NIST.SP.800-53r5.

\bibitem{nist-207} Rose S., Borchert O., Mitchell S., Connelly S. Zero Trust Architecture. NIST SP 800-207. Gaithersburg, MD: NIST, 2020. DOI: 10.6028/NIST.SP.800-207.

\bibitem{nist-40} Souppaya M., Scarfone K. Guide to Enterprise Patch Management Planning. NIST SP 800-40, Rev.~4. Gaithersburg, MD: NIST, 2022. DOI: 10.6028/NIST.SP.800-40r4.

\bibitem{nist-218} Souppaya M., Scarfone K., Dodson D. Secure Software Development Framework (SSDF) Version 1.1. NIST SP 800-218. Gaithersburg, MD: NIST, 2022. DOI: 10.6028/NIST.SP.800-218.

\bibitem{nist-161} Boyens J. et al. Cybersecurity Supply Chain Risk Management Practices for Systems and Organizations. NIST SP 800-161, Rev.~1. Gaithersburg, MD: NIST, 2022. DOI: 10.6028/NIST.SP.800-161r1.

% --- Бази знань та галузеві настанови ---
\bibitem{mitre-attack} MITRE ATT\&CK\textsuperscript{\textregistered}: Adversarial Tactics, Techniques, and Common Knowledge. The MITRE Corporation. URL: \url{https://attack.mitre.org/}.

\bibitem{mitre-defend} MITRE D3FEND\textsuperscript{\texttrademark}: A Knowledge Graph of Cybersecurity Countermeasures. The MITRE Corporation. URL: \url{https://d3fend.mitre.org/}.

\bibitem{mitre-car} MITRE Cyber Analytics Repository (CAR). The MITRE Corporation. URL: \url{https://car.mitre.org/}.

\bibitem{cve} CVE: Common Vulnerabilities and Exposures. The MITRE Corporation. URL: \url{https://www.cve.org/}.

\bibitem{capec} CAPEC: Common Attack Pattern Enumeration and Classification. The MITRE Corporation. URL: \url{https://capec.mitre.org/}.

\bibitem{yara} Álvarez V. M. YARA: The pattern matching swiss knife for malware researchers. URL: \url{https://virustotal.github.io/yara/}.

\bibitem{sigma} Sigma: Generic Signature Format for SIEM Systems. SigmaHQ. URL: \url{https://github.com/SigmaHQ/sigma}.

\bibitem{stix} STIX Version 2.1: Structured Threat Information Expression. OASIS Standard. URL: \url{https://docs.oasis-open.org/cti/stix/v2.1/stix-v2.1.html}.

\bibitem{lockheed-kc} Hutchins E. M., Cloppert M. J., Amin R. M. Intelligence-Driven Computer Network Defense Informed by Analysis of Adversary Campaigns and Intrusion Kill Chains. \emph{Leading Issues in Information Warfare \& Security Research}. 2011. Vol.~1. P.~80--106.

\bibitem{pyramid} Bianco D. The Pyramid of Pain. Enterprise Detection \& Response, 2013. URL: \url{https://detect-respond.blogspot.com/2013/03/the-pyramid-of-pain.html}.

% --- Право ЄС ---
\bibitem{nis2} Directive (EU) 2022/2555 of the European Parliament and of the Council of 14 December 2022 on measures for a high common level of cybersecurity across the Union (NIS 2 Directive). \emph{Official Journal of the EU}. L 333. P.~80--152. URL: \url{https://eur-lex.europa.eu/eli/dir/2022/2555/oj}.

\bibitem{gdpr} Regulation (EU) 2016/679 (General Data Protection Regulation). \emph{Official Journal of the EU}. L 119. P.~1--88. URL: \url{https://eur-lex.europa.eu/eli/reg/2016/679/oj}.

\bibitem{cra} Regulation (EU) 2024/2847 on horizontal cybersecurity requirements for products with digital elements (Cyber Resilience Act). \emph{Official Journal of the EU}. 2024. URL: \url{https://eur-lex.europa.eu/eli/reg/2024/2847/oj}.

\bibitem{budapest} Convention on Cybercrime (Budapest Convention). ETS No.~185. Budapest: Council of Europe, 2001. URL: \url{https://www.coe.int/en/web/cybercrime/the-budapest-convention}.

\bibitem{enisa} ENISA Threat Landscape 2024. Athens: European Union Agency for Cybersecurity, 2024. URL: \url{https://www.enisa.europa.eu/publications/enisa-threat-landscape-2024}.

% --- Знакові кампанії ЗПК (Україна та світ) ---
\bibitem{notpetya} ESET Research. New TeleBots backdoor: First evidence linking Industroyer to NotPetya. ESET WeLiveSecurity, 2018. URL: \url{https://www.welivesecurity.com/}.

\bibitem{industroyer} Cherepanov A., Lipovsky R. Industroyer: Biggest Threat to Industrial Control Systems since Stuxnet. ESET White Paper. Bratislava: ESET, 2017.

\bibitem{industroyer2} Slowik J., Hultquist J. et al. Industroyer2: Nozomi Networks Labs Analyses the IEC~104 Payload: technical report. 2022. URL: \url{https://cert.gov.ua/article/39518}.

\bibitem{wipers2022} Dobrovolsky G. et al. (CERT-UA, ESET, Microsoft). WhisperGate, HermeticWiper, IsaacWiper, CaddyWiper and the 2022 destructive campaigns against Ukraine: joint analyses. 2022.

\bibitem{stuxnet} Falliere N., Murchu L. O., Chien E. W32.Stuxnet Dossier. Version~1.4. Symantec Security Response, 2011.

% --- Монографії та підручники ---
\bibitem{sikorski} Sikorski M., Honig A. Practical Malware Analysis: The Hands-On Guide to Dissecting Malicious Software. San Francisco: No Starch Press, 2012. 800~p.

\bibitem{szor} Szor P. The Art of Computer Virus Research and Defense. Boston: Addison-Wesley, 2005. 744~p.

\bibitem{aycock} Aycock J. Computer Viruses and Malware. New York: Springer, 2006. 227~p. (Advances in Information Security, Vol.~22).

\bibitem{anderson} Anderson R. Security Engineering: A Guide to Building Dependable Distributed Systems. 3rd ed. Indianapolis: Wiley, 2020. 1232~p. URL: \url{https://www.cl.cam.ac.uk/~rja14/book.html}.

\bibitem{saltzer} Saltzer J. H., Schroeder M. D. The Protection of Information in Computer Systems. \emph{Proceedings of the IEEE}. 1975. Vol.~63, No.~9. P.~1278--1308. DOI: 10.1109/PROC.1975.9939.

\bibitem{buryachok} Бурячок В. Л., Толубко В. Б., Хорошко В. О., Толюпа С. В. Інформаційна та кібербезпека: соціотехнічний аспект: підручник / за заг. ред. В. Б. Толубка. Київ: ДУТ, 2015. 288~с.

\bibitem{csecurity} Кібербезпека: освіта, наука, техніка: електронне фахове наукове видання / Київський столичний університет імені Бориса Грінченка. ISSN 2663-4023. URL: \url{https://csecurity.kubg.edu.ua/}.
